% Shape: structured proof in Lamport's hierarchical style, kept to three levels.
% Evidence for each rule: ../../mathematical-writing/ (keys in references.md there).
% Compiles with pdflatex. Packages: amsmath, amssymb, amsthm; the steps use the
% kernel's enumerate, so no structured-proof package is needed.
\documentclass{article}
\usepackage{amsmath,amssymb,amsthm}

\newtheorem{theorem}{Theorem}
\newtheorem{lemma}[theorem]{Lemma}
\theoremstyle{definition}
\newtheorem{definition}[theorem]{Definition}

% Step names <level>position, as in Lamport 2012: the reference to a step is its own
% name only, so a level-2 step is cited as <2>1, never as 3.<2>1.
\makeatletter
\renewcommand{\p@enumii}{}
\renewcommand{\p@enumiii}{}
\makeatother
\newenvironment{steps}{%
  \renewcommand{\theenumi}{\ensuremath{\langle 1\rangle}\arabic{enumi}}%
  \renewcommand{\labelenumi}{\theenumi.}%
  \renewcommand{\theenumii}{\ensuremath{\langle 2\rangle}\arabic{enumii}}%
  \renewcommand{\labelenumii}{\theenumii.}%
  \renewcommand{\theenumiii}{\ensuremath{\langle 3\rangle}\arabic{enumiii}}%
  \renewcommand{\labelenumiii}{\theenumiii.}%
  \setlength{\leftmargini}{3.2em}\setlength{\leftmarginii}{3.2em}\setlength{\leftmarginiii}{3.2em}%
  \begin{enumerate}}{\end{enumerate}}
% Three levels only: enumerate allows four, this shape stops at three on purpose.

\newcommand{\sketch}[1]{\par\noindent\textit{Proof sketch.} #1\par\smallskip}
\newcommand{\uses}[1]{\par\noindent{\small\textsc{Uses:} #1.}\par\smallskip}
\newcommand{\pf}{\textit{Proof:}\ }
\newcommand{\qedstep}{\textsc{Q.e.d.}}

\begin{document}

\section*{The skeleton}

\begin{theorem}\label{thm:skeleton}
One sentence with every hypothesis inside it.
\end{theorem}

\begin{proof}
\uses{Definition 2.1, Lemma 3.2}
\sketch{One line that gives the idea, for example: induct on $n$; the step reduces to Lemma 3.2.}
\begin{steps}
\item\label{s:first} First claim. \pf by Definition 2.1.
\item\label{s:second} Second claim.
  \begin{enumerate}
  \item\label{s:second:a} First sub-claim. \pf by \ref{s:first}.
  \item \qedstep\ \pf by \ref{s:second:a} and Lemma 3.2.
  \end{enumerate}
\item \qedstep\ \pf by \ref{s:first} and \ref{s:second}.
\end{steps}
\end{proof}

\section*{A calculation as the lowest level}

% A chain of relations with one hint per link (Dijkstra 2000); on paper it works
% best as the lowest level of a proof (Lamport 2012).
\begin{align*}
  & a + 1 \\
  \leq\ & \{\ \text{hypothesis 1}\ \} \\
  & b^2 - 3 \\
  =\ & \{\ \text{definition of } c\ \} \\
  & c - 4
\end{align*}

\section*{Worked example}

\begin{theorem}\label{thm:sqrt2}
No rational number $r$ satisfies $r^2 = 2$.
\end{theorem}

\begin{proof}
\sketch{Write $r$ in lowest terms; both numerator and denominator turn out even.}
\begin{steps}
\item\label{e:suffices} It suffices to assume integers $p, q$ with $q \neq 0$,
  $\gcd(p, q) = 1$ and $p^2 = 2q^2$, and derive a contradiction.
  \pf every rational is $p/q$ in lowest terms, and $(p/q)^2 = 2$ exactly when $p^2 = 2q^2$.
\item\label{e:peven} $p$ is even.
  \pf $p^2 = 2q^2$ is even by \ref{e:suffices}, and the square of an odd integer is odd.
\item\label{e:qeven} $q$ is even.
  \begin{enumerate}
  \item\label{e:k} $p = 2k$ for some integer $k$. \pf by \ref{e:peven}.
  \item\label{e:qsq} $q^2 = 2k^2$. \pf substitute \ref{e:k} into $p^2 = 2q^2$ from
    \ref{e:suffices} and divide by 2.
  \item \qedstep\ \pf $q^2$ is even by \ref{e:qsq}, so $q$ is even, as in \ref{e:peven}.
  \end{enumerate}
\item \qedstep\ \pf by \ref{e:peven} and \ref{e:qeven}, 2 divides both $p$ and $q$,
  which contradicts \ref{e:suffices}.
\end{steps}
\end{proof}

\section*{The same result as a paragraph proof}

% Keep a paragraph when the proof is short or one idea; still give the idea first
% and name each warrant (Lamport 2012, Fuller 2014).
\begin{proof}[Paragraph proof of Theorem~\ref{thm:sqrt2}]
The idea is that lowest terms cannot survive the equation. Write $r = p/q$ with
$\gcd(p, q) = 1$, so that $p^2 = 2q^2$. Then $p^2$ is even, hence $p$ is even, say
$p = 2k$. Substituting gives $q^2 = 2k^2$, so $q$ is even too, which contradicts
$\gcd(p, q) = 1$.
\end{proof}

\end{document}
